% !TEX program = lualatex
\documentclass[11pt,a4paper]{ltjsarticle}
\usepackage{amsmath,amssymb,amsthm,mathtools}
\usepackage[colorlinks=true,urlcolor=blue,linkcolor=blue]{hyperref}
\usepackage{luatexja-fontspec}
\setmainjfont{HaranoAjiMincho-Regular}
\newtheorem{theorem}{定理}[section]
\newtheorem{lemma}[theorem]{補題}
\newtheorem{proposition}[theorem]{命題}
\newtheorem{remark}[theorem]{注意}
\renewcommand{\proofname}{証明}
\newcommand{\PP}{\mathbb P}
\newcommand{\CC}{\mathbb C}
\newcommand{\Gr}{\operatorname{Gr}}
\newcommand{\Pois}{\operatorname{Pois}}
\newcommand{\Sym}{\operatorname{Sym}}
\newcommand{\PGL}{\operatorname{PGL}}
\newcommand{\GL}{\operatorname{GL}}
\newcommand{\SL}{\operatorname{SL}}
\title{Poisson cubic fourfoldのrank4成分のスキーム滑らかさ}\author{研究草稿（著者情報未記入）}
\date{2026年11月12日・v2}
\begin{document}
\maketitle
\begin{abstract}
cubic fourfoldの相対Poissonスキームの階数四成分について、
Jacobi方程式の線形化の非零36次整数小行列式からgeneric smoothnessを証明する。
集合論的成分分類と次元公式はnc-cubic論文Aにあるため再証明しない。
rank2および成分交叉のscheme構造は未解決である。
\end{abstract}
\section{対象と既知の入力}\label{sec:input}
$V$を6次元複素ベクトル空間、$S=\Sym(V^*)$とし、$U\subset\PP(S_3)$を滑らかな三次超曲面の開集合とする。$[F]\in U$に対し$X_F=V(F)\subset\PP(V)$と置く。相対Poissonスキーム$\mathfrak I\to U$の各ファイバーは
\[
 \Pois(X_F)=\{\pi\in H^0(X_F,\wedge^2T_{X_F}):[\pi,\pi]=0\}.
\]
以下、その非零部分$\mathfrak I^\times$を考える。$F$のスカラー変更の際には$P$のスカラーを逆に変更し、$\pi=\iota_{dF}P$を不変にする。この局所的自明化の規約は次元と既約性に影響しない。
正規形と相対Poissonモジュライの集合論的分類はnc-cubic論文Aを入力とし、ここでは再証明しない。\n\section{階数四成分のスキーム構造}
nc-cubic論文Aの
\href{https://github.com/uedakazushi/nc-cubic/blob/main/papers/support/papers/sections/fermat_relative_poisson_components.tex}{相対Poisson分類}
を入力とする。rank4 Poisson三重ベクトル$P_0=e_{123}+e_{456}$の$\GL_6$軌道は20次元開軌道、
固定した$P_0$に対する三次式の空間は射影19次元である。
従ってrank4層は39次元の既約局所閉集合である。
この集合論的分類は論文Aに既に証明があり、本稿では繰り返さない。

\begin{proposition}[階数四成分の一般点でのscheme滑らかさ]\label{prop:rank4-generically-smooth}
nc-cubic論文Aの正規形・成分分類を入力とすると、相対Poissonスキーム$\mathfrak I$は階数四成分の一般点で滑らか、したがってその成分はgeneric reducedである。特に分離Fermat点
\[
 F_0=\sum_{i=1}^6x_i^3,\qquad P_0=e_{123}+e_{456}
\]
で相対Poissonスキームは滑らかである。
\end{proposition}
\begin{proof}
$P=P_0$を固定して$F$の変形を考える。次数三の混合単項式（両方の三変数ブロックを含むもの）は36個である。$F_0$におけるJacobi方程式の線形化をこれら36係数に制限し、20個の三重添字ごとに混合三次単項式36個の係数を取ると、整数係数の$720\times36$行列$M$を得る。付属の \texttt{computations/verify\_poisson\_cubic\_fourfold\_tangent.py} は、零から厳密に$M$を構築する。列の順序は同コード内の三次単項式の五重ループの辞書式順序とし、行の順序は三重添字の辞書式順序、その内部で同じ単項式順序とする。

行番号（零始まり）
\[
\begin{split}
&36,38,45,47,57,59,72,81,93,153,155,165,167,189,\\
&201,258,259,260,270,271,272,295,296,307,308,332,\\
&344,381,383,417,474,475,476,511,512,548
\end{split}
\]
で選んだ正方小行列の行列式は
\[
 \det M_{\rm sub}=-39346408075296537575424\ne0
\]
である。よって$F$方向のJacobi線形化の階数は少なくとも36。分離型の二つの純三次成分は20次元の解族を与えるので階数は高々36、従って正確に36である。

$P$の安定開軌道が20次元、$[F]$の全空間が55次元なので、相対ambientの次元は75。Jacobi方程式の接空間の次元は高々$75-36=39$。一方、この既知の39次元既約性によりその点を通る既約なrank四stratumの次元は39。従って局所環のKrull次元と埋込次元はともに39であり、点は滑らかである。滑らかさの開性とrank四stratumの既約性からその一般点も滑らかとなる。
\end{proof}

\section{未解決問題と検証範囲}
位相的分類はnc-cubic論文Aを正本とする。未解決事項は(a)その独立監査、(b)rank二成分のgeneric reducedness・交叉・embedded associated components、(c)商の境界、(d)新規性の徹底的な先行研究照合である。これらを証明済みとはしない。
\texttt{computations/verify\_poisson\_cubic\_fourfold\_tangent.py}に非零小行列式の厳密検算を収録。今回は再実行・実組版していない。
\begin{thebibliography}{9}
\bibitem{Pym} B.~Pym, \emph{Constructions and classifications of projective Poisson varieties}, Lett. Math. Phys. \textbf{108} (2018), 573--632, \url{https://doi.org/10.1007/s11005-017-0984-5}.
\bibitem{Mayanskiy} E.~Mayanskiy, \emph{Poisson cohomology of two Fano threefolds}, J. Algebra \textbf{424} (2015), 21--45, \url{https://arxiv.org/abs/1303.5916}. 三次元の比較であり四次元の分類として引用しない。
\bibitem{MM} H.~Matsumura and P.~Monsky, \emph{On the automorphisms of hypersurfaces}, J. Math. Kyoto Univ. \textbf{3} (1964), 347--361. 書誌と適用条件は最終稿で再照合する。
\bibitem{Internal} \emph{Poisson normal form for smooth cubic hypersurfaces}, 現行研究原稿 \url{https://github.com/uedakazushi/nc-cubic/blob/main/papers/support/papers/sections/poisson_geometry.tex}. 既発表文献ではない。
\end{thebibliography}
\end{document}
